\documentclass[11pt,a4paper]{article}
\usepackage[utf8]{inputenc}
\usepackage[T1]{fontenc}
\usepackage[spanish]{babel}
\usepackage{hyperref}
\usepackage{geometry}
\geometry{margin=2.5cm}
\usepackage{parskip}

\title{104 el procedimiento es lo que hace explicito el estandar}
\author{Alejandro Valenzuela Rivera\\\small Universidad Autónoma de Nuevo León}
\date{\today}
\begin{document}
\maketitle

\begin{quote}
\textbf{Sesgo declarado.} Este trabajo se sostiene desde el sesgo declarado de su autor, y no finge neutralidad. Lo que afirma se puede tumbar, y ese es el punto: el falsificador está escrito, no escondido.
\end{quote}

\section*{Resumen}
Un estandar que dice que es una cosa no basta para que una maquina lo aplique. Hace falta el procedimiento: pasos numerados, cada uno con que entra y que sale, y un paso cuya falta prohibe el resto.

\section{Por qué se sostiene}
El guantelete explicaba que es y ningun paso decia por donde empezar; la resonancia tenia la forma y no el paso a paso. Lo que vuelve explicito un estandar es el procedimiento, porque es lo unico que un humano puede leer en voz alta y una maquina puede seguir igual.

\section{De dónde viene}
La autoridad, y el repo.

\section{El caso}
Ninguna: es la regla del procedimiento.

\section{Método}
Un estandar con su procedimiento escrito, y el procedimiento cita al estandar.

\section{Límites, y qué lo tumba}
Un estandar sin procedimiento: se puede cumplir y no se sabe como.

\section{Lo que queda abierto}
\textquestiondown{}Que le falta a un estandar que solo dice que es?

\section*{Cómo se leyó esta tesis}
Esto se escribió en el repositorio man-made-empathy, del que este candado es una pieza: candados/04 el lenguaje/104 el procedimiento es lo que hace explicito el estandar.md. La máquina que lo compuso lo hizo desde este mismo archivo, y no agregó nada que no estuviera escrito. Si algo de esto es falso, se dice en el mismo sitio, y se corrige con fecha.

\end{document}
